\section*{अध्याय \ref{ch:g1:the-five-senses} --- पाँच इंद्रियाँ}

\begin{solution}{exo:g1:the-five-senses:1}
दृष्टि --- \omterm{def:g1:the-five-senses:sense}{आँखें}; श्रवण --- \omterm{def:g1:the-five-senses:sense}{कान}; घ्राण --- नाक; स्वाद --- जीभ;
स्पर्श --- \omterm{def:g1:the-five-senses:sense}{त्वचा}।
\end{solution}

\begin{solution}{exo:g1:the-five-senses:2}
(क) \omterm{def:g1:the-five-senses:sense}{आँखें}। (ख) \omterm{def:g1:the-five-senses:sense}{त्वचा}। (ग) नाक। (घ) \omterm{def:g1:the-five-senses:sense}{कान}।
\end{solution}

\begin{solution}{exo:g1:the-five-senses:3}
\omterm{def:g1:the-five-senses:sense}{त्वचा}, यानी स्पर्श का अंग। वह \omterm{def:g1:my-body:parts}{सिर} से पाँव तक पूरे शरीर को ढके रहती
है।
\end{solution}

\begin{solution}{exo:g1:the-five-senses:4}
हाथों की \omterm{def:g1:the-five-senses:sense}{त्वचा} ने संदेश भेजा (ठंडी, सख़्त, नुकीली); मस्तिष्क ने उसे
समझा और चीज़ का नाम बताया।
\end{solution}

\begin{solution}{exo:g1:the-five-senses:5}
स्वाद और घ्राण। स्वाद का आधा काम नाक करती है, इसीलिए नाक बंद होने पर
खाना फीका लगता है।
\end{solution}

\begin{solution}{exo:g1:the-five-senses:6}
जैसे: सूरज को कभी सीधे मत देखो --- \omterm{def:g1:the-five-senses:sense}{आँखों} की हिफ़ाज़त; ध्वनियाँ धीमी
रखो --- \omterm{def:g1:the-five-senses:sense}{कानों} की; गरम खाने पर फूँक मारो --- जीभ की; \omterm{def:g1:the-five-senses:sense}{कान} या नाक में
कुछ मत ठूँसो --- \omterm{def:g1:the-five-senses:sense}{कानों} और नाक की; पढ़ने के लिए अच्छा उजाला ---
\omterm{def:g1:the-five-senses:sense}{आँखों} की।
\end{solution}

\begin{solution}{exo:g1:the-five-senses:7}
क्योंकि बहुत तेज़ ध्वनि \omterm{def:g1:the-five-senses:sense}{कानों} को नुक़सान पहुँचाती है, और चोट खाया \omterm{def:g1:the-five-senses:sense}{कान}
हमेशा के लिए ख़राब रह सकता है। कान-ढक्कन मशीनों के शोर को इतना धीमा
रखते हैं कि उम्र भर सुनाई देता रहे।
\end{solution}

\begin{solution}{exo:g1:the-five-senses:8}
उत्तर खुला है --- ज़्यादातर लोग तीन में से दो या तीन के नाम बता देते
हैं। \omterm{def:g1:the-five-senses:sense}{आँखों} की कोई मदद न होने पर भी नाक जानी-पहचानी चीज़ों को अच्छी तरह
पहचान लेती है।
\end{solution}

\begin{solution}{exo:g1:the-five-senses:9}
\omterm{def:g1:the-five-senses:sense}{आँख} सिर्फ़ \emph{देखती} है --- वह प्रकाश जुटाती है और एक संदेश भेजती
है। समझने का काम मस्तिष्क करता है: वह संदेश पाता है और कहता है ``यह
दादी हैं'', ``यह मेरा स्कूल है''। इंद्रियाँ जुटाती हैं; मस्तिष्क
समझता है।
\end{solution}

\begin{solution}{exo:g1:the-five-senses:10}
पढ़ना: उँगलियों के पोरों से --- स्पर्श दृष्टि की जगह ले लेता है, और
पन्ने पर उभरे बिंदुओं को टटोलता है। सड़क पार करना: \omterm{def:g1:the-five-senses:sense}{कानों} से --- आती
और जाती गाड़ियों को सुनकर --- और अकसर एक छड़ी के साथ, जिसकी नोक से
स्पर्श आगे की ज़मीन जाँचता रहता है।
\end{solution}
